\chapter{Traduction automatique --- NLLB-600M}

\section{Identification}

\begin{center}
\begin{tabular}{ll}
\toprule
Modèle         & \texttt{facebook/nllb-200-distilled-600M} \\
Beam search    & $B = 5$ faisceaux \\
Statut scoring & \textbf{Implémenté} \\
Fichier        & \texttt{DataPipelines/TranslationModel.py} \\
\bottomrule
\end{tabular}
\end{center}

\section{Benchmark FLORES+ (200 phrases, 8 paires de langues)}

\begin{center}
\begin{tabular}{lrrrrr}
\toprule
\textbf{Modèle} & FR$\to$EN & EN$\to$FR & FR$\to$AR & AR$\to$FR & ms/phrase \\
\midrule
\textbf{NLLB-600M} & 41.0 & \textbf{46.3} & 13.9 & \textbf{29.7} & \textbf{130--180} \\
Qwen3.5-4B & \textbf{41.2} & 44.7 & 14.5 & 29.3 & 660--900 \\
TGemma-4B  & 37.9 & 40.4 & \textbf{14.5} & 26.4 & 620--940 \\
\bottomrule
\end{tabular}
\end{center}

NLLB retenu : gagne 5--6/8 directions, 4--5$\times$ plus rapide.
Résultats complets dans \texttt{scoring\_rayan/\_archive/}.

\section{Fonction de scoring de confiance}

\noindent\textit{Implémentation dans \texttt{TranslationModel.compute\_confidence()}.}

\medskip
Soit $K$ le nombre de chunks, $\ell_k$ le log-prob normalisé best-beam,
$p_k = \exp(\ell_k)$, et $r_k = \text{output\_len}_k / \text{input\_len}_k$
le ratio de longueur.
%definir P(barre)

\begin{equation}
  \phi = 1 - \frac{|\{k : r_k < 0.3 \;\text{ou}\; r_k > 4\}|}{K} \cdot 0.5
  \label{eq:trad_penalty}
\end{equation}

\begin{equation}
  s_{\text{trad}}(x, \hat{y})
  = \bigl(0.7 \cdot \bar{p} + 0.3 \cdot \min_k p_k\bigr) \cdot \phi
  \;\in (0, 1]
  \label{eq:trad}
\end{equation}

\paragraph{Seuils.}

\begin{center}
\begin{tabular}{lll}
\toprule
$s \geq 0.70$ & Traduction fiable & Exploitable \\
$0.40 \leq s < 0.70$ & Traduction incertaine & Vérification recommandée \\
$s < 0.40$ & Très incertaine & Relecture humaine \\
\bottomrule
\end{tabular}
\end{center}

\textbf{Métadonnées loguées :} \texttt{num\_chunks}, \texttt{src\_lang},
\texttt{tgt\_lang}, \texttt{mean\_beam\_prob}, \texttt{min\_beam\_prob},
\texttt{outlier\_chunks}.

\section{Justification scientifique}

NLLB~\cite{nllb2022} ; calibration beam NMT~\cite{ott2018analyzing} ;
Quality Estimation non supervisée~\cite{fomicheva2020unsupervised}.

\section{Validation empirique}

Aucun dataset d'évaluation annoté n'est disponible à ce jour. Le
\texttt{compute\_confidence()} est actuellement validé par inspection manuelle
sur des cas réels et des scénarios limites (\textit{smoke testing}).

\noindent\textbf{Métriques à mesurer lors de la constitution du dataset :}
\begin{itemize}
  \item Corrélation entre le score de confiance et la métrique de qualité
    appropriée (ex. BLEU, chrF++) ;
  \item Calibration (ECE) : un score de confiance de 0.8 correspond-il à
    environ 80\,\% de réussite sur le dataset ?
  \item Pouvoir discriminant (AUC-ROC) : le score sépare-t-il les traductions
    correctes des incorrectes ?
\end{itemize}

\noindent\textbf{Statut :} \textit{À implémenter.}


% # Composite: mean weighted 0.7, min weighted 0.3 (penalises worst-case chunks)
%        score = (0.7 * mean_beam_prob + 0.3 * min_beam_prob) * outlier_penalty
% la formule de calcul est aussi composite a étudié